\section{Conclusion}
FedGrid combines solar-capacity grouping, robust filtering, clipping and reference-based weighting for federated battery control. In the five-seed Ausgrid archive, the final configuration, which changes filter order and clipping strength, reduces observed mean annual Normalized-attack damage from approximately \$88 to \$11 per household. ACT retraining shows a reduction in the same direction, subject to preprocessing overlap. The exact five-seed signed-rank comparisons do not establish a difference at the 0.05 level, and reward poisoning remains a weakness.

The evaluation separates economic protection from energy utilisation and admission behaviour. Clean FedAvg substantially reduces curtailment relative to the rule, while bill differences are modest and battery throughput increases. Filter admission does not ensure positive aggregation weight, just as a bound on aggregate norm does not guarantee a favourable direction or safe controller behaviour.

These findings support evaluating federated energy control through attack-specific bill damage, electrical outcomes and realised contribution weights. Partition-isolated preprocessing, independent evaluation, matched information assumptions and household-level cost accounting would strengthen the evidence connecting aggregation decisions to reliable and economically useful collaboration.
